\chapter{System Architecture and Protocol Design}

This chapter describes the overall system architecture of the LE-PSI implementation
and presents the full two-round protocol as realised in software.

\section{Protocol Overview}

The LE-PSI protocol derived from DKLLMR23~\cite{dottling2023efficient} proceeds in
two rounds plus a local computation phase:

\begin{enumerate}
    \item \textbf{Server Digest (Round 1, Server $\to$ Client):}
          The server $\mathcal{P}_1$ builds a Merkle tree over its dataset $S$ and
          publishes the root $\mathsf{pp}$ as a public parameter.
          $|\mathsf{pp}| = O(D \log n)$ ring elements.
    \item \textbf{Client Encryption (Round 2, Client $\to$ Server):}
          For each $c_j \in C$, the client $\mathcal{P}_2$ computes a Ring-LWE
          ciphertext $\mathsf{ct}_j = \mathsf{LE.Enc}(\mathsf{pp}, c_j, \mathsf{msg})$
          and sends $\{\mathsf{ct}_j\}_{j=1}^{m}$ to the server.
    \item \textbf{Intersection Detection (Server, local):}
          For each $s_i \in S$, the server generates a witness $w_i$ from the
          Merkle tree and attempts decryption of every $\mathsf{ct}_j$. If
          $c_j = s_i$, decryption succeeds; this signals $s_i \in S \cap C$.
\end{enumerate}

\begin{figure}[H]
\centering
\renewcommand{\arraystretch}{1.3}
\begin{tabular}{p{5.5cm} c p{5.5cm}}
\hline
\multicolumn{1}{c}{\textbf{Server} ($S$, $n$ records)} & &
\multicolumn{1}{c}{\textbf{Client} ($C$, $m$ records)} \\
\hline
\vspace{2mm}
$(\mathsf{pp}, \{sk_i, pk_i\}) \leftarrow \mathsf{ServerInit}(S)$ & & \\
& $\xrightarrow{\;\;\mathsf{pp}\;\;}$ & \\
& &
$\{\mathsf{ct}_j\} \leftarrow \mathsf{ClientEnc}(\mathsf{pp}, C)$ \\
& $\xleftarrow{\;\;\{\mathsf{ct}_j\}\;\;}$ & \\
$S \cap C \leftarrow \mathsf{Intersect}(\{sk_i\},\{w_i\},\{\mathsf{ct}_j\})$ & & \\
\hline
\end{tabular}
\caption{Two-round LE-PSI protocol}
\label{fig:protocol}
\end{figure}

\section{Software Package Architecture}

The implementation is written in Go~1.24.1 and structured as follows:

\begin{verbatim}
LE-PSI/
+-- pkg/
|   +-- LE/           # Laconic Encryption core
|   |   +-- keys.go   # KeyGen: pk = A*sk + e  (RLWE sample)
|   |   +-- tree.go   # Merkle tree (SQLite + in-memory)
|   |   +-- wit.go    # WitGenMemory: Merkle witness from RAM
|   +-- psi/
|   |   +-- server.go      # ServerInitialize (monolithic, small n)
|   |   +-- client.go      # ClientEncrypt (parallel)
|   |   +-- helpers.go     # CalculateOptimalWorkers
|   |   +-- parameters.go  # SetupLEParameters (D=256 / D=2048)
|   +-- matrix/       # Polynomial ring R_q arithmetic (NTT, GInvMNTT)
+-- internal/storage/ # SQLite Merkle tree I/O
+-- utils/            # Hashing, type conversion
+-- scalability_tests/# HPC benchmark harness (Phase II)
+-- cmd/Flare/        # Sanctions screening CLI
\end{verbatim}


\section{Dual-Mode Security Configuration}

The system supports two operating security levels selected via the environment variable
\texttt{PSI\_SECURITY\_LEVEL}:

\begin{table}[H]
\centering
\caption{Security mode parameters}
\label{tab:modes}
\begin{tabular}{@{}lcccc@{}}
\toprule
\textbf{Mode} & \textbf{$D$} & \textbf{Quantum security} & \textbf{Mem/record} & \textbf{Use case} \\
\midrule
Default ($D$=256)  & 256  & $\approx 64$-bit  & $\sim$35\,MB  & Large-scale benchmarking \\
128-bit PQ         & 2048 & $\approx 128$-bit & $\sim$280\,MB & High-security deployment \\
\bottomrule
\end{tabular}
\end{table}

The \texttt{SetupLEParameters()} function selects ring dimension, modulus, and tree
depth based on this environment variable.

\section{Data Flow}

\begin{enumerate}
    \item \textbf{Data ingestion:} Server records are loaded from a financial
          transaction SQLite database (6.36M records, PaySim-derived dataset).
          Records are hashed to $\mathbb{Z}_{2^L}$ leaf indices using SHA-256.
    \item \textbf{Key generation:} For each record, a Ring-LWE key pair
          $(pk_i, sk_i)$ is generated. Public keys are inserted into the Merkle
          tree via \texttt{LE.Upd()}.
    \item \textbf{Tree loading:} After construction, the entire tree is loaded
          from SQLite into a \texttt{MemoryTree} data structure.
    \item \textbf{Client encryption:} Client records are hashed, and one
          Ring-LWE ciphertext is generated per client record.
    \item \textbf{Batched intersection:} Server records are processed in batches
          of size $B$. For each batch: witnesses are generated, decryption
          checks are performed, and witnesses are freed.
    \item \textbf{Output:} Matching server record hashes are collected and
          reported.
\end{enumerate}

\clearpage
